\summarychapter{1}{Away with the King!}{The Jesus Movement}

\begin{inbrief}
Christianity is the only major religion whose central event is the humiliation of its God. The chapter
follows Jesus from the divided Palestine of his day, through his preaching of the kingdom of God and
his clash with the Jewish parties, to the cross, and argues that the company he gathered was the
beginning of the church.
\question{Did Jesus have anything to do with founding the Christian church, and if so, how did he give
it its special character?}
\end{inbrief}

The chapter opens at the cross. Crucifixion was a barbarous death kept for agitators, pirates and
slaves; Jewish law cursed it and Cicero wanted the very word kept from Roman ears. The notice over
Jesus read ``Jesus of Nazareth, King of the Jews'' (INRI). Pilate meant it as a last insult to the
Jews, but, like the cross itself, it took on a special meaning for the followers of Jesus.

\ssection{Jesus and the Church}
Jesus was a Jew in family, scripture and religion, and many have doubted that he meant to found a
church. Albert Schweitzer saw a man possessed by the coming end of the world, Rudolf Bultmann a prophet
calling for a radical decision; others speak only of an invisible brotherhood of love. Against this
anti-institutional view Shelley points to the gospels, which show Jesus retracing Israel's steps---Egypt, the Jordan, the wilderness, twelve apostles like twelve tribes, teaching like Moses, mighty
deeds---and being the faithful Son where Israel failed. He gathered disciples, taught them the life
of the kingdom of God and bound them in a new covenant of forgiveness and love. That company lacked the
laws, officials and creeds of later Christendom, but it was a society apart that learned to refuse
rival loyalties. In that sense Jesus founded the church.

\ssection{Palestine in Jesus' Day}
Palestine was a crossroads: two million or more people under Rome, divided by region, religion and
politics. The Jews, about half the population, hated their rulers and the Greek (Hellenistic) culture
they brought, and differed only about how to resist. Popular faith, fed by the prophets and by Daniel,
expected a messiah, a final struggle between God and Satan, the resurrection of the dead, a last
judgment and God's everlasting kingdom.

Four parties read the crisis differently. The \emph{Pharisees} (``separated ones'') kept every detail
of the law, avoided the unclean and were respected for piety and patriotism. The \emph{Sadducees}
represented the wealthy priestly families who ran the temple and the Sanhedrin, liked Greco-Roman ways
and profited from Roman rule. The \emph{Zealots}, inspired by the Maccabees, wanted armed revolt and
kept guerrilla bands in Galilee. The \emph{Essenes} withdrew from a temple they thought corrupt to
monastic communities in the wilderness; scholars link them with Qumran and the Dead Sea Scrolls. Jesus
had to win followers without being confused with any of them.

\begin{sidebox}{Judaism Now and Then}
The destruction of the temple in AD 70 changed Judaism. Rome's crushing of revolt ended the Zealot and
Essene protests, and the Sadducees' religion vanished with the temple. The Pharisees, the ``book
people'' who ordered all of life by the Hebrew Scriptures, survived and became today's Rabbinic Judaism.
\end{sidebox}

\ssection{Jesus' Ministry}
Jesus began with John the Baptist, who preached repentance and baptism at the ford of the Jordan and
called himself only the forerunner of the Messiah. After his own baptism Jesus started his mission in
Galilee as a traveling teacher in the village synagogues. His theme was the kingdom of God, which
Shelley takes as both future and present: not a territory but the rule of a personal, gracious God,
already active in Jesus himself. His healings were signs of the coming age, though he told the healed
to keep silent so as not to be taken for a magician. Crowds gathered, and the Pharisees, who resented a
teacher not trained by their scribes, challenged his credentials.

\ssection{Jesus' Message}
The challenge let Jesus contrast grace with self-righteousness. In his parable of two men at prayer,
the Pharisee boasts of fasting and tithing and despises others; the tax collector, compromised by his
work for Rome, only begs for mercy, and he is the one justified. The Pharisees' piety rested on keeping
hundreds of laws; the Jesus movement rested on trust in God's mercy. Jesus chose twelve apostles
(``sent ones''), a motley group from fishing boats and tax tables, and taught them that his kingdom
has another kind of greatness, governed by love and service. His popularity peaked a year before his
arrest, when he fed over five thousand pilgrims and many wanted to make him king; he withdrew. Wary of
messianic titles the crowds would misread, he described his role in the images of Isaiah's Suffering
Servant and Zechariah's humble king.

\ssection{The Last Week}
Five days before the Passover of his death Jesus entered Jerusalem riding a donkey to shouts of ``Hosanna,''
the one occasion on which he openly claimed the role of Messiah, and the next day he drove the traders
from the temple. Talk of a messiah frightened the Sadducees, whose position depended on Rome and who
dreaded a revolt; the Pharisees saw a breaker of the Sabbath who undermined the law. The rival parties
joined to have him killed and, fearing a riot, used Judas Iscariot, one of the Twelve, who agreed for
thirty pieces of silver---nearly four months' wages---to lead them to him secretly.

\ssection{The New Covenant}
At the Passover meal Jesus gave his disciples bread and a cup and spoke of ``the new covenant in my
blood.'' Behind the words stood the Exodus and Sinai, but above all Jeremiah's promise of a covenant
written on the heart, with sins forgiven: through his death a new people of God became possible. The
puzzled disciples would understand within weeks. Afterward, in Gethsemane on the Mount of Olives, Jesus
prayed that the cup might pass yet accepted his Father's will; Judas came with an armed crowd, and
Jesus was taken to the high priest Caiaphas.

\ssection{The Trial and Death}
The Sanhedrin brushed justice aside, found two witnesses and condemned Jesus for blasphemy. Lacking
power to execute, it took him at dawn to the Roman governor, Pontius Pilate, at the Antonia fortress.
Pilate saw a petty religious dispute, but he feared a charge of sheltering a traitor to Caesar more
than a festival riot, and gave Jesus to the soldiers. At Golgotha, outside the city, he was nailed to
the cross under the sign naming him king of the Jews; after hours of agony, watched by a few despairing
followers, he cried ``It is finished'' and died. Joseph of Arimathea laid him in a rock tomb in his
garden and closed it with a heavy stone---for Shelley, an odd but telling opening for the story of the church.

\begin{keyfacts}
\keypeople{Jesus of Nazareth; John the Baptist; the twelve apostles; Judas Iscariot; Caiaphas; Pontius
Pilate; Joseph of Arimathea. Interpreters cited: Albert Schweitzer, Rudolf Bultmann.}
\keydates{%
  \kd{AD 70}{Destruction of the temple; only the Pharisees' form of Judaism survives}
  \kd{1946}{Discovery of the Dead Sea Scrolls (date as given in the book)}}
\keyterms{Kingdom of God; Messiah; Pharisees, Sadducees, Zealots, Essenes; Sanhedrin; apostles; new
covenant; Suffering Servant.}
\end{keyfacts}
